\section{数学准备与算法描述}
\subsection{矩阵范数与精确极因子}
对 $W\in\R^{m\times n}$，记 $r=\min(m,n)$。本文使用 Frobenius 内积
$\inner{A}{B}=\operatorname{tr}(A^\top B)$，并区分三个范数：
\begin{equation}
 \Frob{A}=\Bigl(\sum_i s_i(A)^2\Bigr)^{1/2},\quad
 \op{A}=s_1(A),\quad \nuc{A}=\sum_i s_i(A).
\end{equation}
由 Cauchy--Schwarz 不等式，$\Frob{A}\leq\nuc{A}\leq\sqrt r\Frob{A}$。后一个关系中的 $\sqrt r$ 是本报告随机收敛常数出现维度依赖的来源，不能在隐藏常数时忽略。

若 $M=U_q\Sigma_qV_q^\top$ 为仅包含正奇异值的紧致 SVD，定义
\begin{equation}\label{eq:polar}
 O^\star=\polar(M)=U_qV_q^\top;\qquad \polar(0)=0.
\end{equation}
当 $M$ 满列秩或满行秩时，$O^\star$ 是半正交矩阵。秩亏时，$O^\star$ 是对应支撑子空间上的部分等距映射；一般不能写 $(O^\star)^\top O^\star=I$。这一约定与保留零奇异值对应基向量的任意半正交补全不同，但均可成为谱范数单位球上的线性子问题解。紧致约定能保证零动量不产生更新，也与多项式迭代对零奇异值的处理一致。

\begin{lemma}[对偶关系与方向性质]\label{lem:duality}
对任意矩阵 $M$，式\eqref{eq:polar}中的 $O^\star$ 满足
\begin{equation}\label{eq:dual}
 \op{O^\star}\leq1,\qquad \Frob{O^\star}\leq\sqrt r,\qquad
 \inner{M}{O^\star}=\nuc M=\max_{\op D\leq1}\inner M D.
\end{equation}
\end{lemma}
\begin{proof}
由紧致 SVD，$O^\star$ 的非零奇异值均为 1，数量为 $q\leq r$，故前两个结论成立。迹的循环性质给出 $\inner{M}{O^\star}=\operatorname{tr}(\Sigma_q)=\nuc M$。另一方面，von Neumann 迹不等式给出
$\inner M D\leq\sum_i s_i(M)s_i(D)\leq\nuc M\op D$，故任何可行 $D$ 的目标值都不超过 $\nuc M$。零矩阵情形直接成立。
\end{proof}

因此，$-O^\star$ 解决谱范数单位球上的线性最小化子问题，即
\begin{equation}
 \arg\min_{\|D\|_2\leq1}\inner{M}{D}\ \ni\ -O^\star.
\end{equation}
负号不能省略；$O^\star$ 本身最大化该内积。还需区分\emph{归一化方向}与\emph{局部二次模型的完整最优步}：
\begin{equation}\label{eq:quadraticmodel}
 \arg\min_D\left\{\inner M D+\frac{1}{2\gamma}\op D^2\right\}
 \ \ni\ -\gamma\nuc M\,O^\star.
\end{equation}
这是因为固定 $\op D=s$ 时，最小线性项为 $-s\nuc M$，再对 $s\geq0$ 最小化即可。Muon 使用的是 $-\eta O^\star$；不能把它与式\eqref{eq:quadraticmodel}在同一个数值步长下混为一谈。

\subsection{用于证明的精确版本}
记 $G_t$ 为在当前参数 $W_t$ 处取得的批量随机梯度，$\alpha=1-\beta$。本文理论版本如下。
\begin{algorithm}[htbp]
\caption{精确极因子 Muon：本文的理论分析对象}\label{alg:exact}
\KwIn{$W_1$，迭代预算 $T$，步长 $\eta>0$，动量参数 $0\leq\beta<1$，批量 $B$}
在 $W_1$ 处采样 $G_1$，初始化 $M_1\leftarrow G_1$\;
\For{$t=1,\ldots,T$}{
 \If{$t\geq2$}{在 $W_t$ 处采样 $G_t$\;
 $M_t\leftarrow\beta M_{t-1}+(1-\beta)G_t$\;}
 $O_t^\star\leftarrow\polar(M_t)$，零矩阵按式\eqref{eq:polar}处理\;
 $W_{t+1}\leftarrow W_t-\eta O_t^\star$\;
}
\KwOut{参数序列；用于理论评价时，均匀抽取一个 $W_t$}
\end{algorithm}

这里采用 $M_1=G_1$，便于明确初始误差；它与常见实现中从零缓冲区开始的约定不同。附录中的确定性二次问题与随机矩阵实验遵循算法\ref{alg:exact}。该版本没有偏差修正、权重衰减、梯度裁剪或 Nesterov 外推。

\subsection{有限步实现与参数分组}
实际实现先形成 EMA 动量 $M_t=\beta M_{t-1}+(1-\beta)G_t$，再可选地使用
\begin{equation}\label{eq:nesterov}
 N_t=\beta M_t+(1-\beta)G_t.
\end{equation}
该式采用参考代码的 EMA 插值约定。令 $X_0=N_t/(\Frob{N_t}+10^{-7})$，对较长维度作转置以降低矩阵乘法成本，再执行
\begin{equation}\label{eq:ns}
 X_{j+1}=aX_j+(bX_jX_j^\top+c(X_jX_j^\top)^2)X_j,
 \quad (a,b,c)=(3.4445,-4.7750,2.0315).
\end{equation}
最后转回原形状。本文五步配置采用 $j=0,\ldots,4$；历史 CPU 语言模型与新增 CUDA 图像实验均使用 float32，参考实现的 bfloat16 版本与此存在精度差异。若 $X_0=U\Sigma V^\top$，一次迭代保留奇异向量，并将奇异值映射为 $\phi(s)=as+bs^3+cs^5$。

有限步输出不必等于 $UV^\top$。经验系数的目标是用少量矩阵乘法获得有效方向，并不保证对任意谱，迭代越多就越接近精确极因子。$\phi(0)=0$ 也说明多项式本身不会补出缺失的秩。

应用更新为
\begin{equation}\label{eq:practical}
 W_{t+1}=(1-\eta_t\lambda)W_t-\eta_t s_{m,n}O_t,
 \qquad s_{m,n}=\sqrt{\max(1,m/n)}.
\end{equation}
语言模型只把 Transformer 块内的二维线性层权重分给 Muon。词嵌入、位置嵌入、输出层权重、偏置及归一化参数采用 AdamW。精确 SVD 对照也保留式\eqref{eq:nesterov}、形状缩放及相同参数分组，仅替换正交化方式。

新增 ResNet-18 实验将隐藏卷积核展平为“输出通道 $\times$ 其余维度”的二维矩阵。输入 stem、分类 head 和一维参数仍由 AdamW 更新。$K=0$ 消融定义为仅保留前置 Frobenius 归一化、不执行多项式；它不是零更新，也不等于未经归一化的 SGD 动量。
\begin{algorithm}[htbp]
\caption{实际混合 Muon 的参数分组与更新}\label{alg:hybrid}
\KwIn{隐藏矩阵组 $\mathcal H$，辅助组 $\mathcal A$，$K$，动量 $\beta$，步长与衰减}
初始化各隐藏参数的动量缓冲区为零\;
\For{每个训练批量}{
 在同一组当前参数处完成前向、反向，取得全部梯度\;
 \For{$W\in\mathcal H$}{
 将卷积梯度按输出通道展平，$M\leftarrow\beta M+(1-\beta)G$\;
 $N\leftarrow\beta M+(1-\beta)G$，$O\leftarrow\mathrm{NS}_K(N)$\;
 $W\leftarrow(1-\eta\lambda)W-\eta\sqrt{\max(1,m/n)}\,\mathrm{reshape}(O)$\;
 }
 对 $\mathcal A$ 执行一次 AdamW 更新，偏置及归一化参数不衰减\;
}
\end{algorithm}

对稠密 $m\times n$ 矩阵，薄 SVD 的典型算术量为 $O(mnr)$；转置后使用较小 Gram 矩阵的 $J$ 步多项式也为 $O(Jmnr)$，但具体常数、并行方式及数值精度不同。AdamW 的逐坐标更新为 $O(mn)$，通常保存两个矩估计；Muon 的矩阵参数部分仅持久保存一个动量矩阵，但正交化仍需临时工作空间。因此，状态数量较少与单步更新更快是两个不同的问题。


\FloatBarrier
